% !TeX root = ../../progress-sem1.tex
% สารบัญคำศัพท์ทางเทคนิคและคำย่อ (glossary). Shared by every report.
%: Add your own terms and definitions below, following the example format.
\newcommand{\tableraglf}[1]{\RaggedRight{#1}}
\listofvocab
\begin{flushleft}
  \begin{xltabular}{\textwidth}{@{}>{\RaggedRight\arraybackslash}p{0.31\textwidth}@{\quad=\quad}>{\RaggedRight\arraybackslash}X@{}}
    UX/UI Design / การออกแบบประสบการณ์ผู้ใช้และส่วนต่อประสานผู้ใช้ & \tableraglf{กระบวนการออกแบบที่ครอบคลุมสองส่วนหลัก คือ UX (User Experience) ซึ่งเน้นการออกแบบภาพรวมของประสบการณ์ที่ผู้ใช้ได้รับขณะใช้งานผลิตภัณฑ์หรือบริการ เช่น ความง่ายในการใช้งาน ความสะดวก และความพึงพอใจ และ UI (User Interface) ซึ่งเน้นการออกแบบส่วนต่อประสานที่ผู้ใช้มองเห็นและโต้ตอบด้วยโดยตรง เช่น ปุ่มกด เมนู สี และรูปแบบตัวอักษร โดยทั้งสองส่วนทำงานร่วมกันเพื่อให้ผลิตภัณฑ์ดิจิทัล เช่น เว็บไซต์หรือแอปพลิเคชัน ใช้งานได้ง่าย มีประสิทธิภาพ และตอบสนองความต้องการของผู้ใช้ได้อย่างแท้จริง} \\

    UX/UI development life cycle (UX/UI DLC) & \tableraglf{วงจรการพัฒนางานออกแบบ UX/UI ซึ่งเป็นกระบวนการที่มีขั้นตอนต่อเนื่องกันตั้งแต่การเริ่มต้นจนถึงการนำไปใช้งานจริง โดยทั่วไปประกอบด้วยขั้นตอนหลัก เช่น การวิจัยและทำความเข้าใจผู้ใช้ (Research) การกำหนดปัญหาและความต้องการ (Define) การออกแบบแนวคิดและโครงร่าง (Ideate) การสร้างต้นแบบ (Prototype) และการทดสอบกับผู้ใช้จริง (Testing) ทั้งนี้ในแต่ละรอบของวงจรมักมีการนำผลตอบรับกลับมาปรับปรุงซ้ำ (Iteration) เพื่อให้ผลิตภัณฑ์ตอบโจทย์ผู้ใช้ได้ดียิ่งขึ้น} \\

    Usability Testing / การทดสอบการใช้งาน & \tableraglf{กระบวนการทดสอบการใช้งาน โดยให้ผู้เข้าร่วมทดสอบลองใช้งานระบบและหาจุดที่ใช้งานยากหรือติดขัด ซึ่งในขณะที่ทำการทดสอบ ผู้เข้าร่วมการทดสอบจำเป็นต้องคิดไปด้วยพูดไปด้วย (Thinking-Aloud) โดยนักออกแบบจะคอยสังเกตพฤติกรรมการใช้งานและรับฟังความเห็นของผู้เข้าร่วมทดสอบขณะทำภารกิจที่กำหนด} \\

    Comparative Usability Testing & \tableraglf{วิธีการทดสอบการใช้งานเชิงคุณภาพกับผู้เข้าร่วมทดสอบ โดยให้ผู้เข้าร่วมทดสอบทดลองใช้งานต้นแบบผลิตภัณฑ์สองรูปแบบที่มีเป้าประสงค์เดียวกันเพื่อเปรียบเทียบประสิทธิภาพการใช้งานและให้ข้อคิดเห็นหรือข้อเสนอแนะเกี่ยวกับรูปแบบหรือผลิตภัณฑ์นั้น ๆ พร้อมทั้งเลือกรูปแบบใดรูปแบบหนึ่งที่เห็นว่าเหมาะสมกับเป้าประสงค์นั้น ๆ} \\

    Contextual Feedback / ข้อเสนอแนะเชิงบริบท & \tableraglf{ข้อคิดเห็นหรือคำแนะนำที่กล่าวเฉพาะเจาะจงกับสิ่งใดสิ่งหนึ่ง เพื่อให้คำแนะนำมีความเกี่ยวข้องและนำไปใช้ได้ตรงจุดมากขึ้น} \\

    Wireframe / โครงร่างหน้าจอ & \tableraglf{โครงร่างพื้นฐานแสดงตำแหน่งองค์ประกอบบนหน้าจอ โดยยังไม่เน้นสีสันหรือรายละเอียดด้านภาพ ใช้สื่อสารการจัดวางในขั้นต้นของการออกแบบ} \\

    Prototype / ต้นแบบ & \tableraglf{ต้นแบบผลิตภัณฑ์ที่สร้างขึ้นเพื่อทดสอบแนวคิดก่อนพัฒนาจริง สามารถโต้ตอบได้ในระดับหนึ่ง แบ่งเป็นสามระดับตามความน้อย-มากของรายละเอียดในต้นแบบ คือ low, mid, high-fidelity prototype} \\

    Design System & \tableraglf{ชุดมาตรฐานการออกแบบที่รวบรวมองค์ประกอบ กฎเกณฑ์ และแนวทางปฏิบัติไว้ในที่เดียว เช่น สี ตัวอักษร ไอคอน และ component ต่าง ๆ เพื่อให้การออกแบบผลิตภัณฑ์มีความสอดคล้องและเป็นมาตรฐานเดียวกันทั้งระบบ} \\

    Insight / ข้อมูลเชิงลึก & \tableraglf{ความเข้าใจที่ได้จากการวิเคราะห์ข้อมูลหรือพฤติกรรมผู้ใช้ ซึ่งเปิดเผยปัญหาหรือความต้องการที่ซ่อนอยู่} \\

    User / ผู้ใช้ & \tableraglf{บุคคลที่มีปฏิสัมพันธ์โดยตรงกับผลิตภัณฑ์ ระบบ หรือบริการ เพื่อบรรลุเป้าหมายของตนเอง} \\

    Generative AI / ปัญญาประดิษฐ์แบบรู้สร้าง & \tableraglf{เทคโนโลยีปัญญาประดิษฐ์ที่สามารถสร้างเนื้อหาใหม่ เช่น ข้อความ ภาพ เสียง หรือโค้ดจากคำสั่ง (prompt)} \\

    Design Intent / ความต้องการทางการออกแบบ & \tableraglf{เจตนาหรือเป้าหมายเชิงการออกแบบที่นักออกแบบต้องการสื่อสาร ใช้เป็นแนวทางตัดสินใจในการออกแบบ} \\

    Component & \tableraglf{ส่วนประกอบย่อยที่นำกลับมาใช้ซ้ำได้ใน UI เช่น ปุ่มหรือช่องกรอกข้อมูล ช่วยให้การออกแบบสอดคล้องกัน} \\

    Two alternative UI & \tableraglf{UI ที่ใช้สำหรับการทำ comparative usability testing โดย UI ทั้งสองแบบต้องถูกทำขึ้นเพื่อเป้าประสงค์เดียวกัน แต่มีความแตกต่างกันในด้านการใช้งาน เช่น microcopy, layout} \\

    Human-in-the-loop & \tableraglf{แนวคิดที่มนุษย์เข้ามาใช้วิจารณญาณตรวจสอบและปรับปรุงผลลัพธ์จาก Generative AI เพื่อความถูกต้องเหมาะสม} \\

    Stakeholder / ผู้มีส่วนได้เสีย & \tableraglf{บุคคลหรือกลุ่มที่มีส่วนได้ส่วนเสียกับโครงการ เช่น ผู้บริหารหรือผู้ว่าจ้าง ซึ่งมีบทบาทในการตัดสินใจหรือสนับสนุนโครงการ} \\

   % Cognitive Overload & \tableraglf{ภาวะที่ได้รับข้อมูลมากเกินความสามารถในการประมวลผลของสมอง ส่งผลให้สับสนหรือตัดสินใจผิดพลาด} \\

    Requirement / ความต้องการ & \tableraglf{เงื่อนไขที่ผลิตภัณฑ์ต้องมีเพื่อตอบสนองความต้องการผู้ใช้หรือเป้าประสงค์ของโครงการ} \\

    Usability Tester / ผู้เข้าร่วมทดสอบ & \tableraglf{บุคคลที่เข้าร่วมทดสอบการใช้งานผลิตภัณฑ์ พร้อมให้ข้อคิดเห็นเพื่อนำไปปรับปรุงการออกแบบ} \\

    Persona Artifact & \tableraglf{เอกสารจำลองลักษณะผู้ใช้กลุ่มเป้าหมาย เช่น ชื่อ อายุ พฤติกรรม และเป้าหมาย เพื่อให้ทีมออกแบบเข้าใจผู้ใช้ชัดเจนขึ้น} \\
    %  & \tableraglf{ } \\

    %  & \tableraglf{ } \\

    %  & \tableraglf{ } \\

    %  & \tableraglf{ } \\

    %  & \tableraglf{ } \\
  \end{xltabular}
\end{flushleft}
